\section{Conclusiones}
\label{sec:conclusiones}

Este trabajo presentó el diseño, implementación y evaluación empírica del Sistema de Gestión de
Pre-Sustentaciones UTEQ, siguiendo Design Science Research \cite{peffers2007} de forma instanciada y
verificable. Los 90 requisitos del SRS v1.0.1 (64 funcionales, 26 no funcionales) están implementados o planificados
según su estado declarado (36 de 50 Must verificados con prueba automatizada real, 72\,\%) y el flujo
académico completo se verificó funcionando end-to-end contra una topología de producción real. La evaluación empírica combina
estadística descriptiva e inferencial genuina --- el hallazgo más sólido es la reducción de latencia de
13.9$\times$ por efecto de caché Redis, con significancia estadística fuerte ($p < 10^{-10}$) y tamaño
de efecto máximo --- y auditoría de seguridad multi-herramienta que confirma ausencia de inyección SQL en
el código propio; se corrigió la única vulnerabilidad de dependencias que afectaba código servido al
navegador (\texttt{@angular/core}, detectada por ZAP), y quedan abiertas 14 vulnerabilidades en
herramientas de build que no se envían a producción. A esa auditoría automatizada se sumó una revisión
manual dirigida del modelo de autorización que encontró y cerró 6 fallos reales de control de acceso
(IDOR de lectura y de escritura, un endpoint de eliminación sin permiso administrativo, y
mass-assignment en la creación de usuarios) que las herramientas automatizadas no detectan por no
evaluar lógica de negocio por recurso (Sección~\ref{sec:seguridad-corregida}).

La contribución más distintiva de este informe no es técnica sino metodológica: el proceso de
elaboración de este mismo proyecto incluyó la detección real de datos académicos fabricados en fases
anteriores (usabilidad, rendimiento web, evidencia de trazabilidad, integridad de configuración de
despliegue, custodia de consentimientos informados) y su corrección sistemática, documentada de forma
trazable en \texttt{docs/mediciones/DATA-PROVENANCE.md} y en la bitácora de cada fase del proyecto. Esa
corrección no fue cosmética: implicó reportar una cobertura de código de 63.17\,\% de líneas en ese
momento (hoy 70.00\,\% global, ver Sección~\ref{sec:evaluacion}) en vez de un ``$>$60\,\%'' previamente afirmado
sin evidencia, y declarar la usabilidad como ``sin datos'' en vez de mantener un puntaje fabricado (hoy aplicado con
$n=4$ de fecha verificable, en vez de las 15 respuestas recolectadas cuando 11 tienen una fecha que no
se sostiene --- ver \texttt{docs/mediciones/sus/SUS-RESULTS.md}).

Esa misma lógica de corrección continua se aplicó también a la evaluación del propio proyecto: una
revisión de checkpoint durante el desarrollo calificó el estado del proyecto en ese momento con 6.53
sobre 10. En vez de tratar esa cifra como un resultado a superar en el reporte final, se documenta aquí
como el punto de partida real, y las correcciones descritas en este informe (el cierre de los 6 fallos de
autorización de la Sección~\ref{sec:seguridad-corregida}, la sincronización de las secuencias de
PostgreSQL, los índices en columnas de alto tráfico, y el crecimiento de 109 a 559 pruebas automatizadas
con la cobertura subiendo de 38.88\,\% a 70.00\,\% de líneas) como las acciones concretas y verificables
tomadas después de esa evaluación. Este informe no reporta ni infiere una nota final --- eso corresponde
a una evaluación posterior por el docente-director, no a una autoevaluación del propio equipo.

Las brechas que quedan (cobertura de pruebas incompleta, ausencia de datos SUS reales, despliegue público
pendiente, dependencias vulnerables sin corregir) se declaran explícitamente en la
Sección~\ref{sec:trabajo-futuro} en vez de ocultarse o rellenarse con contenido que simule completitud.
Esta decisión --- priorizar un informe más corto pero verificable sobre uno más extenso pero con
contenido fabricado --- es, en sí misma, la conclusión metodológica central de este trabajo: en un
proyecto con antecedentes reales de fabricación de evidencia, la integridad de lo que se reporta importa
más que la completitud aparente de lo que se reporta.
\newpage
